\documentclass[CJK,13pt]{beamer}
\usepackage{CJKutf8}
\usepackage{beamerthemesplit}
\usetheme{Malmoe}
\useoutertheme[footline=authortitle]{miniframes}
\usepackage{amsmath}
\usepackage{amssymb}
\usepackage{graphicx}
\usepackage{eufrak}
\usepackage{color}
\usepackage{slashed}
\usepackage{simplewick}
\usepackage{tikz}
\usepackage{tcolorbox}
\usepackage{ulem}
\graphicspath{{../figures/}}
%%figures
\def\lfig#1#2{\includegraphics[width=#1 in]{#2}}
\def\tfig#1#2{\includegraphics[height=#1 in]{#2}}
\def\addfig#1#2{\begin{center}\includegraphics[width=#1 in]{#2}\end{center}}
\def\question{\includegraphics[width=0.3in]{why.jpg}\,}
\def\answer{\includegraphics[width=0.3in]{baozou_haha.png}\,}
\def\wulian{\includegraphics[width=0.18in]{emoji_wulian.jpg}}
\def\bigwulian{\includegraphics[width=0.35in]{emoji_wulian.jpg}}
\def\bye{\includegraphics[width=0.18in]{emoji_bye.jpg}}
\def\bigbye{\includegraphics[width=0.35in]{emoji_bye.jpg}}
\def\huaixiao{\includegraphics[width=0.18in]{emoji_huaixiao.jpg}}
\def\bighuaixiao{\includegraphics[width=0.35in]{emoji_huaixiao.jpg}}
\def\jianxiao{\includegraphics[width=0.18in]{emoji_jianxiao.jpg}}
\def\bigjianxiao{\includegraphics[width=0.35in]{emoji_jianxiao.jpg}}
\def\haoqi{\includegraphics[width=0.18in]{emoji_haoqi.jpg}}
%% colors
\def\blacktext#1{{\color{black}#1}}
\def\bluetext#1{{\color{blue}#1}}
\def\redtext#1{{\color{red}#1}}
\def\darkbluetext#1{{\color[rgb]{0,0.2,0.6}#1}}
\def\skybluetext#1{{\color[rgb]{0.2,0.7,1.}#1}}
\def\cyantext#1{{\color[rgb]{0.,0.5,0.5}#1}}
\def\greentext#1{{\color[rgb]{0,0.7,0.1}#1}}
\def\darkgray{\color[rgb]{0.2,0.2,0.2}}
\def\lightgray{\color[rgb]{0.6,0.6,0.6}}
\def\gray{\color[rgb]{0.4,0.4,0.4}}
\def\blue{\color{blue}}
\def\red{\color{red}}
\def\orange{\color[rgb]{1.,0.8,0.}}
\def\green{\color{green}}
\def\darkgreen{\color[rgb]{0,0.4,0.1}}
\def\darkblue{\color[rgb]{0,0.2,0.6}}
\def\skyblue{\color[rgb]{0.2,0.7,1.}}
%%control
\def\diag{\mathrm{diag}\,}
\def\heaviside{\,\mathrm{h}\,}
\def\bral{\left(\begin{array}{l}}
\def\brar{\end{array}\right)}
\def\brall{\left(\begin{array}{ll}}
\def\brarr{\end{array}\right)}
\def\bralll{\left(\begin{array}{lll}}
\def\brarrr{\end{array}\right)}
\def\branchl{\left\{\begin{array}{l}}
\def\branchr{\end{array}\right.}
\def\branchll{\left\{\begin{array}{ll}}
\def\branchrr{\end{array}\right.}
\def\branchlll{\left\{\begin{array}{lll}}
\def\branchrrr{\end{array}\right.}
\def\sfgamma#1{\,\Gamma\left( #1 \right)\,}
\def\be{\begin{equation}}
\def\ee{\nonumber\end{equation}}
\def\bea{\begin{eqnarray}}
\def\eea{\nonumber\end{eqnarray}}
\def\bch{\begin{CJK}{UTF8}{gbsn}}
\def\ech{\end{CJK}}
\def\bitem{\begin{itemize}}
\def\eitem{\end{itemize}}
\def\bcenter{\begin{center}}
\def\ecenter{\end{center}}
\def\bex{\begin{minipage}{0.2\textwidth}\includegraphics[width=0.6in]{jugelizi.png}\end{minipage}\begin{minipage}{0.76\textwidth}}
\def\eex{\end{minipage}}
\def\chtitle#1{\frametitle{\bch#1\ech}}
\def\bmat#1{\left(\begin{array}{#1}}
\def\emat{\end{array}\right)}
\def\bcase#1{\left\{\begin{array}{#1}}
\def\ecase{\end{array}\right.}
\def\bmini#1{\begin{minipage}{#1\textwidth}}
\def\emini{\end{minipage}}
\def\tbox#1{\begin{tcolorbox}#1\end{tcolorbox}}
\def\pfrac#1#2#3{\left(\frac{\partial #1}{\partial #2}\right)_{#3}}
\def\res#1#2{\,\mathrm{res}\,#1\left(#2\right)\,}
\def\newt#1#2{\left(\begin{array}{c}#1\\ #2\end{array}\right)}
\def\reof#1{\,\mathrm{Re}{\left(#1\right)\,}}
\def\imof#1{\,\mathrm{Im}{\left(#1\right)\,}}
\def\Arg#1{\,\mathrm{Arg}\,#1\,}
%%symbols
\def\intfull{\int_{-\infty}^\infty}
\def\inthalf{\int_0^\infty}
\def\bropt{\,(\ \ \ )}
\def\sech{\mathrm{sech}\,}
\def\csch{\mathrm{csch}\,}
\def\asinh{\mathrm{asinh}\,}
\def\acosh{\mathrm{acosh}\,}
\def\sone{$\star$}
\def\stwo{$\star\star$}
\def\sthree{$\star\star\star$}
\def\sfour{$\star\star\star\star$}
\def\sfive{$\star\star\star\star\star$}
\def\rint{{\int_\leftrightarrow}}
\def\roint{{\oint_\leftrightarrow}}
\def\stdHf{{\textit{\r H}_f}}
\def\deltaH{{\Delta \textit{\r H}}}
\def\ii{{\dot{\imath}}}
\def\skipline{{\vskip0.1in}}
\def\skiplines{{\vskip0.2in}}
\def\lagr{{\mathcal{L}}}
\def\hamil{{\mathcal{H}}}
\def\vecT{{\mathbf{T}}}
\def\vecN{{\mathbf{N}}}
\def\vecB{{\mathbf{B}}}
\def\vecv{{\mathbf{v}}}
\def\vecr{{\mathbf{r}}}
\def\vecf{{\mathbf{f}}}
\def\vecg{{\mathbf{g}}}
\def\vecupsilon{{\mathbf{\upsilon}}}
\def\vecu{{\mathbf{u}}}
\def\vecj{{\mathbf{j}}}
\def\vecx{{\mathbf{x}}}
\def\vecy{{\mathbf{y}}}
\def\vecz{{\mathbf{z}}}
\def\veck{{\mathbf{k}}}
\def\vecp{{\mathbf{p}}}
\def\vecn{{\mathbf{n}}}
\def\vecA{{\mathbf{A}}}
\def\vecP{{\mathbf{P}}}
\def\vecsigma{{\mathbf{\sigma}}}
\def\hatJn{{\hat{J_\vecn}}}
\def\hatJx{{\hat{J_x}}}
\def\hatJy{{\hat{J_y}}}
\def\hatJz{{\hat{J_z}}}
\def\hatj#1{\hat{J_{#1}}}
\def\hatphi{{\hat{\phi}}}
\def\hatq{{\hat{q}}}
\def\hatpi{{\hat{\pi}}}
\def\vel{\upsilon}
\def\Dint{{\mathcal{D}}}
\def\adag{{\hat{a}^\dagger}}
\def\bdag{{\hat{b}^\dagger}}
\def\cdag{{\hat{c}^\dagger}}
\def\ddag{{\hat{d}^\dagger}}
\def\hata{{\hat{a}}}
\def\hatb{{\hat{b}}}
\def\hatc{{\hat{c}}}
\def\hatd{{\hat{d}}}
\def\hatD{{\,\hat{D}}}
\def\hatN{{\hat{N}}}
\def\hatH{{\hat{H}}}
\def\hatp{{\hat{p}}}
\def\Fup{{F^{\mu\nu}}}
\def\Fdown{{F_{\mu\nu}}}
\def\newl{\nonumber \\}
\def\vece{\mathrm{e}}
\def\calM{{\mathcal{M}}}
\def\calT{{\mathcal{T}}}
\def\calR{{\mathcal{R}}}
\def\barpsi{\bar{\psi}}
\def\baru{\bar{u}}
\def\barv{\bar{\upsilon}}
\def\qeq{\stackrel{?}{=}}
\def\ftf{\stackrel{\mathcal{FT}}{\Longrightarrow}}
\def\ftb{\stackrel{\mathcal{FT}}{\Longleftarrow}}
\def\ftfb{\stackrel{\mathcal{FT}}{\Longleftrightarrow}}
\def\ltf{\stackrel{\mathcal{LT}}{\Longrightarrow}}
\def\ltb{\stackrel{\mathcal{LT}}{\Longleftarrow}}
\def\ltfb{\stackrel{\mathcal{LT}}{\Longleftrightarrow}}
\def\torder#1{\mathcal{T}\left(#1\right)}
\def\rorder#1{\mathcal{R}\left(#1\right)}
\def\contr#1#2{\contraction{}{#1}{}{#2}#1#2}
\def\trof#1{\mathrm{Tr}\left(#1\right)}
\def\trace{\mathrm{Tr}}
\def\comm#1{\ \ \ \left(\mathrm{used}\ #1\right)}
\def\tcomm#1{\ \ \ (\text{#1})}
\def\slp{\slashed{p}}
\def\slk{\slashed{k}}
\def\calp{{\mathfrak{p}}}
\def\veccalp{\mathbf{\mathfrak{p}}}
\def\Tthree{T_{\tiny \textcircled{3}}}
\def\pthree{p_{\tiny \textcircled{3}}}
\def\dbar{{\,\mathchar'26\mkern-12mu d}}
\def\erf{\mathrm{erf}}
\def\const{\mathrm{const.}}
\def\pheat{\pfrac p{\ln T}V}
\def\vheat{\pfrac V{\ln T}p}
%%units
\def\fdeg{{^\circ \mathrm{F}}}
\def\cdeg{^\circ \mathrm{C}}
\def\atm{\,\mathrm{atm}}
\def\angstrom{\,\text{\AA}}
\def\SIL{\,\mathrm{L}}
\def\SIkm{\,\mathrm{km}}
\def\SIyr{\,\mathrm{yr}}
\def\SIGyr{\,\mathrm{Gyr}}
\def\SIV{\,\mathrm{V}}
\def\SImV{\,\mathrm{mV}}
\def\SIeV{\,\mathrm{eV}}
\def\SIkeV{\,\mathrm{keV}}
\def\SIMeV{\,\mathrm{MeV}}
\def\SIGeV{\,\mathrm{GeV}}
\def\SIcal{\,\mathrm{cal}}
\def\SIkcal{\,\mathrm{kcal}}
\def\SImol{\,\mathrm{mol}}
\def\SIN{\,\mathrm{N}}
\def\SIHz{\,\mathrm{Hz}}
\def\SIm{\,\mathrm{m}}
\def\SIcm{\,\mathrm{cm}}
\def\SIfm{\,\mathrm{fm}}
\def\SImm{\,\mathrm{mm}}
\def\SInm{\,\mathrm{nm}}
\def\SImum{\,\mathrm{\mu m}}
\def\SIJ{\,\mathrm{J}}
\def\SIW{\,\mathrm{W}}
\def\SIkJ{\,\mathrm{kJ}}
\def\SIs{\,\mathrm{s}}
\def\SIkg{\,\mathrm{kg}}
\def\SIg{\,\mathrm{g}}
\def\SIK{\,\mathrm{K}}
\def\SImmHg{\,\mathrm{mmHg}}
\def\SIPa{\,\mathrm{Pa}}
%page
\def\secpage#1#2{\begin{frame}\bcenter{\bf \Huge #1} \ecenter \skipline \tbox{\bcenter #2 \ecenter}\end{frame}}
\def\append#1#2{\secpage{附录#1}{#2}}
\def\thinka#1{\begin{frame}\frametitle{思考题}\bcenter\lfig{0.4}{think0.jpg}  \skipline #1 \ecenter\end{frame}}
\def\thinkb#1{\begin{frame}\frametitle{思考题}\bcenter\lfig{0.5}{think1.jpg}  \skipline #1 \ecenter\end{frame}}
\def\thinkc#1{\begin{frame}\frametitle{思考题}\bcenter\lfig{0.5}{think2.jpg}  \skipline #1 \ecenter\end{frame}}
\def\thinkd#1{\begin{frame}\frametitle{思考题}\bcenter\lfig{0.5}{think3.jpg}  \skipline #1 \ecenter\end{frame}}
\def\thinke#1{\begin{frame}\frametitle{思考题}\bcenter\lfig{0.5}{think4.jpg}  \skipline #1 \ecenter\end{frame}}
\def\thinkf#1{\begin{frame}\frametitle{思考题}\bcenter\lfig{0.5}{think5.jpg}  \skipline #1 \ecenter\end{frame}}
\def\schw{ds^2 = \left(1-\frac{2GM}{r}\right)dt^2 - \left(1-\frac{2GM}{r}\right)^{-1} dr^2 - r^2\left(d\theta^2 + \sin^2\theta d\phi^2\right)}

\input{titlepage.tex}
  \date{}
  \begin{document}
  \bch
  \tpage{15}{Test Particle in Schwarzschild Metric}

  \begin{frame}
    这一讲我们讨论 \sout{地球为什么绕着太阳转} Schwarzschild 度规下的测试粒子的运动。

    \addfig{2}{solar.jpg}
    
  \end{frame}

  \begin{frame}
    \frametitle{守恒动量}
    \addfig{0.5}{think0.jpg}

    设度规的协变形式不依赖于某个坐标分量 $x^\alpha$，即 $g_{\mu\nu,\alpha}=0$ 处处成立，证明沿测地线运动的自由粒子的四维动量的协变分量 $p_\alpha$ 守恒。

  \end{frame}


  \secpage{史瓦西度规里的测试粒子的基本方程}{$$\theta=\frac{\pi}{2}$$
    $$p_t=E$$
    $$p_\phi = -L$$
    $$p^\mu p_\mu = m^2$$
  }

  \begin{frame}
    考虑在史瓦西度规
    $$ ds^2 = \left(1-\frac{2GM}{r}\right)dt^2 - \left(1-\frac{2GM}{r}\right)^{-1} dr^2 - r^2\left(d\theta^2 + \sin^2\theta d\phi^2\right)$$
    中运动的测试粒子。记其四维动量为 $p^\mu$。
  \end{frame}
    
  \begin{frame}
    取粒子运动所在的三维平面为
   {\blue \begin{equation}
      \theta  = \frac{\pi}{2} \label{eq:theta}
    \end{equation}}
    度规
    $$ \diag{\left(1-\frac{2GM}{r}, -\left(1-\frac{2GM}{r}\right)^{-1},-r^2, -r^2\sin^2\theta\right)}$$
    和 $t, \phi$ 无关，所以四维动量的协变分量 $p_t$, $p_\phi$ 均守恒，它们在牛顿力学里都有个响亮的名号： $p_t$ 是能量（这里比牛顿力学多包含了静止质量）， $p_\phi$ 是角动量（多个负号）。我们记
    {\blue
    \begin{equation}
      p_t = E \label{eq:t}
    \end{equation}

    \begin{equation}
      p_\phi = -L \label{eq:phi}
    \end{equation}
    }
    最后，按照静止质量的定义
{\blue    \begin{equation}
      p^\mu p_\mu=m^2 \label{eq:m2}
    \end{equation}}
  \end{frame}
  

  \begin{frame}
    前面的四个方程已经可以完备地描述粒子的运动。

    \addfig{2}{yali.jpg}
    
    就是那么简单，甚至不用算联络！
  \end{frame}

    \secpage{静质量非零粒子的运动}{$$\frac{d^2u}{d\phi^2} + u =q^2+3u^2 $$}
  
  \begin{frame}
    对静止质量 $m>0$ 的粒子，我们研究的基本对象是变量
    $$ u \equiv \frac{GM}{r}$$
    随着方位角 $\phi$ 的变化。
   
    {\scriptsize 这个神奇的变量 $u$ 是从牛顿引力理论里学来的，别问，再问就拿苹果砸你……}
    
    \skipline
    
    Eq.~\eqref{eq:t} 可以写成
    \begin{equation}
      \left(1- 2u\right)\frac{dt}{ds} =  \varepsilon \label{eq:ts}
    \end{equation}
    这里的 $\varepsilon \equiv \frac{E}{m}$ 是单位质量对应的能量。在自然单位制下，它是无量纲量。
  \end{frame}

  \begin{frame}
    Eq.~\eqref{eq:phi} 可以写成
    \begin{equation}
      r^2\frac{d\phi}{ds} = \ell \label{eq:phis}
    \end{equation}
    这里的  $\ell \equiv \frac{L}{m}$ 是单位质量对应的角动量，它具有长度量纲。由于 $GM$ 也是长度量纲，我们自然可以引入一个无量纲量：
    \begin{equation}
      q \equiv \frac{GM}{\ell}
    \end{equation}
    对太阳系的行星而言，这个量都远小于 $1$。对于近圆轨道，$q^2$是太阳的 $GM$（约1.5公里）和行星公转轨道半径之比。
  \end{frame}

  \begin{frame}
    利用 Eq.~\eqref{eq:ts}和 Eq.~\eqref{eq:phis}，Eq.~\eqref{eq:m2} 可以写成
    \begin{equation}
      (1-2u) \left(\frac{\varepsilon}{1-2u}\right)^2 - \frac{1}{1-2u}\left(\frac{dr}{ds}\right)^2 - \frac{\ell^2}{r^2} = 1
    \end{equation}
    再次利用 Eq.~\eqref{eq:phis}，把上式中的 $\frac{dr}{ds}$ 替换为 $\frac{dr}{d\phi}$
    \begin{equation}
     \varepsilon^2- \left(\frac{dr}{d\phi}\frac{\ell}{r^2}\right)^2 - (1-2u)\left(1+\frac{\ell^2}{r^2} \right) = 0
    \end{equation}
    上式两边乘以  $q^2$，得到
    \begin{equation}
      \varepsilon^2q^2 - \left(\frac{du}{d\phi}\right)^2 - (1-2u)\left(q^2 + u^2\right) = 0
    \end{equation}
    两边对 $\phi$ 求导，得到
    {\blue
    \begin{equation}
      \frac{d^2u}{d\phi^2} + u =q^2+3u^2
    \end{equation}
    }
  \end{frame}
  
  \begin{frame}
    最后这个方程是推广了的Binet公式：
   \tbox{$$\frac{d^2u}{d\phi^2} + u =q^2+3u^2 $$}
    在牛顿近似下 $|u| \ll 1$，所以我们暂且忽略掉方程右边的 $3u^2$，得到
    \begin{equation}
      u(\phi) \approx u_0(\phi) = q^2 (1+e\cos\phi) \label{eq:sol}
    \end{equation}
    这里 $e$ 是任意常数。扎实的幼儿园基础告诉我们，Eq.~\eqref{eq:sol} 里的 $u_0(\phi)$ 其实就是牛顿行星理论里的圆锥曲线。

    {\scriptsize 因为 $\phi$ 的零点可以随便选取，所以我们忽略了 Eq.~\eqref{eq:sol} 里的 $\phi$ 可能带有的常数相位。}
  \end{frame}

  
  \begin{frame}
    \frametitle{近圆轨道的修正}
    我们来考虑水星进动的问题。这对应离心率 $e$ 很小的情况。

    我们来考虑形如
    $$ u = Aq^2\left[1+e\cos\left(\frac{\phi}{B}\right)\right]$$
    的解，这里的 $A,B$ 均很接近于$1$。代入相对论 Binet 公式，并忽略最高阶的 $q^4e^2$ 小项
    $$ Ae \left(1 -\frac{1}{B^2}-6Aq^2\right)\cos\left(\frac{\phi}{B}\right) = 1 + 3A^2q^2-A$$
    保留最低阶修正
    $$ A \approx 1+3q^2, B\approx 1+3q^2$$
  \end{frame}
  
  \begin{frame}
    \frametitle{水星进动(Mercury Perihelion Shift)}
    于是广义相对论预言的轨道，相对于牛顿理论有两种修正：轨道半长轴被缩短 $1+3q^2$ 倍，$u(\phi)$ 的周期被增大 $1+3q^2$ 倍。

    \skiplines

    对太阳系而言，由于半长轴的缩短和太阳质量是简并的，除非我们有不借助于行星运动的方法测量太阳质量，否则就无法验证这个效应。

    \skiplines
    
    $u(\phi)$ 周期偏离 $2\pi$，即每运行一周多出 $6\pi q^2$，可以通过观测行星近日点方位的变化得到。

    我们来看一下水星：太阳的 $GM$ 大约是$1.5$公里，水星的轨道半径约为$5.6$千万公里，于是 $6\pi q^2\approx 5.0\times 10^{-7}$。水星的公转周期是$0.24$年。每百年由广义相对论修正引起的进动为：$5\times 10^{-7}\times\frac{100}{0.24}=2.1\times 10^{-4}$——这正是广义相对论提出之前天文学家们苦苦思索不得的 “多余的 $43$ 角秒”。
  \end{frame}


  \secpage{零质量粒子的运动}{$$\frac{d^2u}{d\phi^2} + u = 3u^2 $$}

  \begin{frame}
    对零质量粒子，$p^\mu p_\mu = 0$ 可以写成
    \begin{equation}
     \left(p^r\right)^2 = E^2 - \frac{(1-2u)u^2L^2}{G^2M^2} \label{eq:ds2}
    \end{equation}

    \begin{equation}
      \left(\frac{du}{d\phi}\right)^2 = \frac{G^2M^2\left(p^r\right)^2}{L^2} = q^2 - (1-2u)u^2
    \end{equation}
    我们这里定义类和静止质量非零时类似的一个量
    $$q \equiv \frac{GME}{L} $$
    然后两边对 $\phi$ 求导，得到
    {\blue $$\frac{d^2u}{d\phi^2} + u = 3u^2 $$}
  \end{frame}

  \begin{frame}
    所以零质量粒子的相对论版 Binet 公式为：
    \tbox{$$\frac{d^2u}{d\phi^2} + u = 3u^2 $$}
    我们仍然只考虑 $u$ 很小的情况。

    \skipline
     
    先忽略方程右边的 $3u^2$ 得到的解为
    $$ u \approx q \cos \phi $$
    这里的 $|q|\ll 1$.
    
    {\scriptsize 因为 $\phi$ 的零点可以随便选取，所以我们忽略了$\phi$ 可能带有的常数相位。}

    \skipline
    
    这显然是一条直线的方程（忽略时空弯曲的解），且
    $$q \equiv \frac{GM}{d}$$
    这里$d$ 是太阳到光线的垂直距离。
  \end{frame}  


  \begin{frame}
    把零级解代入到修正项 $3u^2$ 里，得到一级近似解
    \be
    u \approx q\cos\phi + \frac{3q^2}{2} - \frac{q^2}{2}\cos(2\phi)
    \ee
    我们感兴趣的是 $u=0$ 对应的两个 $\phi$ 的解。设这两个解为 $\pm\frac{\pi+\epsilon}{2}$ （显然光线的偏折角即为 $\epsilon$）。保留到 $\epsilon$ 的一级小量，有
    $$ -q\frac{\epsilon}{2} + 2q^2 = 0$$
    即偏折角为 $\epsilon = 4q$。
  \end{frame}

  \begin{frame}
    牛顿力学其实对光线如何偏折没有确定性的结论。因为按牛顿引力公式，光子受力为零，加速度 $a=F/m$ 则是一个不确定量。

    \skipline
    
    如果我们换一种观点，把光子看成质量很小但不是零，且以光速运动的（这在牛顿力学范畴内是允许的），轨迹几乎为直线的测试粒子。牛顿力学的 Binet 公式给出的解为

    \begin{equation}
      u = q^2(1+e\cos\phi) \label{eq:bsol}
    \end{equation}
    这里的 $e\gg 1$。按照前面对非零质量粒子的讨论，这里的 $q = \frac{GM}{d}$，$d$ 是太阳到粒子轨迹的垂直距离（对以光速运动的粒子而言，单位质量的角动量为 $\ell = d$)。Eq.~\eqref{eq:bsol} 中令 $\phi = 0$，应该得到 $u=\frac{GM}{d}=q$，即得到 $q = q^2(1+e)$ ，故$e =  \frac{1}{q}-1\approx\frac{1}{q}$.

      \skipline
    
    Eq.~\eqref{eq:bsol} 中令 $u=0$，给出偏折角 $\approx \frac{2}{e}\approx 2q$。是广义相对论预言的一半。
  \end{frame}  
  
  \begin{frame}
    
    对水星进动和光线经过太阳时的弯曲的观测验证，奠定了\sout{老爱} 广义相对论的不可动摇的江湖地位。

    \addfig{1}{einsteinIPhone.jpg}
  \end{frame}  
\ech
\end{document}




  
